% !TeX root = ../../../main.tex
\begin{ChSixUncaptionedFigure}
\begin{ChSixFloatBody}

\centering
\begin{adjustbox}{max width=\linewidth, max totalheight=0.45\textheight}
\begin{tcolorbox}[ChSix/annopanel]

\begin{minipage}[t]{0.49\textwidth}\vspace{0pt}
   % <-- important
\begin{tcolorbox}[ChSix/annobox,title={Mixtral-8x7B-Instruct v0.1 annotation}]
\ChSixAnnotationTableStyle
\begin{tabular}{>{\raggedright\arraybackslash}p{0.78\textwidth} r}
\toprule
\textbf{Event} & \textbf{Time} \\
\midrule
57 years old& 0 \\
male & 0 \\
recently diagnosed with lepromatous leprosy & -60 \\
on treatment (rifampicin/clofazimine/dapsone) & -60 \\
presented to the hospital & 0 \\
complaints of abdominal distension & 0 \\
complaints of constipation & 0 \\
complaints of vomiting & 0 \\
complaints of 10-kg weight loss & 0 \\
evidence of peripheral lymphadenopathy & 0 \\
distended abdomen & 0 \\
positive shifting dullness & 0 \\
CT showed mural thickening of the terminal ileum & 0 \\
CT showed significantly enlarged mesenteric lymph nodes & 0 \\
CT showed mesenteric fat stranding & 0 \\
CT showed intra-abdominal free fluid & 0 \\
showing a dramatic response to treatment & 0 \\
planned for consolidation by autologous bone marrow transplant & 0 \\
re-admitted to the medical ICU for severe sepsis and multiorgan failure & 0 \\
passed away around 6 months after his initial diagnosis with NHL & 0 \\
despite maintaining a remission status & 0 \\
\bottomrule
\end{tabular}
\end{tcolorbox}
\end{minipage}
\hfill
\begin{minipage}[t]{0.49\textwidth}\vspace{0pt}
   % <-- important
\begin{tcolorbox}[ChSix/annobox,title={Llama-3.1-8B-Instruct annotation}]
\ChSixAnnotationTableStyle
\begin{tabular}{>{\raggedright\arraybackslash}p{0.78\textwidth} r}
\toprule
\textbf{Event} & \textbf{Time} \\
\midrule
18 years old & 0 \\
male & 0 \\
admitted to the hospital & 0 \\
lepromatous leprosy & -720 \\
treatment with rifampicin/clofazimine/dapsone & -720 \\
abdominal distension & -720 \\
constipation & -720 \\
vomiting & -720 \\
weight loss & -720 \\
peripheral lymphadenopathy & -720 \\
distended abdomen & -720 \\
positive shifting dullness & -720 \\
mural thickening of the terminal ileum & -720 \\
enlarged mesenteric lymph nodes & -720 \\
mesenteric fat stranding & -720 \\
intra-abdominal free fluid & -720 \\
abdominal granulomatous infection or neoplastic process & -720 \\
autologous bone marrow transplant & -720 \\
recurrent bacteremia & -720 \\
sepsis & -720 \\
multiorgan failure & -720 \\
passed away & -720 \\
\bottomrule
\end{tabular}
\end{tcolorbox}
\end{minipage}

\end{tcolorbox}
\end{adjustbox}

\end{ChSixFloatBody}
\end{ChSixUncaptionedFigure}